\documentclass[10pt,a4paper]{amsart}
\usepackage{amsmath,amssymb,amsthm,graphicx}
% Retained minimal preamble for the Stacks Project chapter flat.tex (More on Flatness).
% Verbatim excerpt of the chapter preamble: Xy-pic setup (native preamble lines 14-18) and the native
% theorem declarations with the mathematical macros (native preamble lines 154-217). Copied mechanically;
% nothing added, removed or redefined. The wrapper class and the amsmath/amssymb/amsthm/graphicx base
% packages are supplied by the assembler, so the class/externaldocument part of the native preamble is omitted.
% For commutative diagrams we use Xy-pic
\usepackage[all]{xy}

% We use 2cell for 2-commutative diagrams.
\xyoption{2cell}
% Theorem environments.
%
\theoremstyle{plain}
\newtheorem{theorem}[subsection]{Theorem}
\newtheorem{proposition}[subsection]{Proposition}
\newtheorem{lemma}[subsection]{Lemma}

\theoremstyle{definition}
\newtheorem{definition}[subsection]{Definition}
\newtheorem{example}[subsection]{Example}
\newtheorem{exercise}[subsection]{Exercise}
\newtheorem{situation}[subsection]{Situation}

\theoremstyle{remark}
\newtheorem{remark}[subsection]{Remark}
\newtheorem{remarks}[subsection]{Remarks}

\numberwithin{equation}{subsection}

% Macros
%
\def\lim{\mathop{\mathrm{lim}}\nolimits}
\def\colim{\mathop{\mathrm{colim}}\nolimits}
\def\Spec{\mathop{\mathrm{Spec}}}
\def\Hom{\mathop{\mathrm{Hom}}\nolimits}
\def\Ext{\mathop{\mathrm{Ext}}\nolimits}
\def\SheafHom{\mathop{\mathcal{H}\!\mathit{om}}\nolimits}
\def\SheafExt{\mathop{\mathcal{E}\!\mathit{xt}}\nolimits}
\def\Sch{\mathit{Sch}}
\def\Mor{\mathop{\mathrm{Mor}}\nolimits}
\def\Ob{\mathop{\mathrm{Ob}}\nolimits}
\def\Sh{\mathop{\mathit{Sh}}\nolimits}
\def\NL{\mathop{N\!L}\nolimits}
\def\CH{\mathop{\mathrm{CH}}\nolimits}
\def\proetale{{pro\text{-}\acute{e}tale}}
\def\etale{{\acute{e}tale}}
\def\QCoh{\mathit{QCoh}}
\def\Ker{\mathop{\mathrm{Ker}}}
\def\Im{\mathop{\mathrm{Im}}}
\def\Coker{\mathop{\mathrm{Coker}}}
\def\Coim{\mathop{\mathrm{Coim}}}

% Boxtimes
%
\DeclareMathSymbol{\boxtimes}{\mathbin}{AMSa}{"02}

%
% Macros for moduli stacks/spaces
%
\def\QCohstack{\mathcal{QC}\!\mathit{oh}}
\def\Cohstack{\mathcal{C}\!\mathit{oh}}
\def\Spacesstack{\mathcal{S}\!\mathit{paces}}
\def\Quotfunctor{\mathrm{Quot}}
\def\Hilbfunctor{\mathrm{Hilb}}
\def\Curvesstack{\mathcal{C}\!\mathit{urves}}
\def\Polarizedstack{\mathcal{P}\!\mathit{olarized}}
\def\Complexesstack{\mathcal{C}\!\mathit{omplexes}}
% \Pic is the operator that assigns to X its picard group, usage \Pic(X)
% \Picardstack_{X/B} denotes the Picard stack of X over B
% \Picardfunctor_{X/B} denotes the Picard functor of X over B
\def\Pic{\mathop{\mathrm{Pic}}\nolimits}
\def\Picardstack{\mathcal{P}\!\mathit{ic}}
\def\Picardfunctor{\mathrm{Pic}}
\def\Deformationcategory{\mathcal{D}\!\mathit{ef}}

% --- appended 2026-10-12 (Chief): the project own definitions of the reference/slogan/history
% display environments, copied mechanically from the shared Stacks preamble lines 8-12
% (\usepackage{verbatim} plus the three \newenvironment lines, which the chapter preamble keeps
% verbatim from the project preamble). Nothing added, removed or redefined: these environments
% carry no mathematical content and the project itself typesets them as comments. Without these
% three lines a counted statement whose own text uses them fails with Environment slogan undefined.
\usepackage{verbatim}
\newenvironment{reference}{\comment}{\endcomment}
\newenvironment{slogan}{\comment}{\endcomment}
\newenvironment{history}{\comment}{\endcomment}
\usepackage{cleveref}
\crefname{originallemma}{original lemma}{original lemmas}
\Crefname{originallemma}{Original lemma}{Original lemmas}
\makeatletter
\newlabel{algebra-lemma-NAK}{{00DV}{source}{}{sourcepointer-1}{}}
\newlabel{algebra-lemma-NAK@cref}{{[originallemma][0][]00DV}{[1][0][]source}{}{}{}}
\makeatother
\crefname{originallemma}{original lemma}{original lemmas}
\Crefname{originallemma}{Original lemma}{Original lemmas}
\makeatletter
\newlabel{algebra-lemma-construct-fp-module-from-stalk}{{05GJ}{source}{}{sourcepointer-2}{}}
\newlabel{algebra-lemma-construct-fp-module-from-stalk@cref}{{[originallemma][0][]05GJ}{[1][0][]source}{}{}{}}
\makeatother
\crefname{originallemma}{original lemma}{original lemmas}
\Crefname{originallemma}{Original lemma}{Original lemmas}
\makeatletter
\newlabel{algebra-lemma-silly}{{00DS}{source}{}{sourcepointer-3}{}}
\newlabel{algebra-lemma-silly@cref}{{[originallemma][0][]00DS}{[1][0][]source}{}{}{}}
\makeatother
\crefname{originallemma}{original lemma}{original lemmas}
\Crefname{originallemma}{Original lemma}{Original lemmas}
\makeatletter
\newlabel{divisors-lemma-finite-ass}{{05AF}{source}{}{sourcepointer-4}{}}
\newlabel{divisors-lemma-finite-ass@cref}{{[originallemma][0][]05AF}{[1][0][]source}{}{}{}}
\makeatother
\crefname{originallemma}{original lemma}{original lemmas}
\Crefname{originallemma}{Original lemma}{Original lemmas}
\makeatletter
\newlabel{divisors-lemma-ses-weakly-ass}{{05AN}{source}{}{sourcepointer-5}{}}
\newlabel{divisors-lemma-ses-weakly-ass@cref}{{[originallemma][0][]05AN}{[1][0][]source}{}{}{}}
\makeatother
\crefname{originallemma}{original lemma}{original lemmas}
\Crefname{originallemma}{Original lemma}{Original lemmas}
\makeatletter
\newlabel{divisors-lemma-weakly-ass-zero}{{05AP}{source}{}{sourcepointer-6}{}}
\newlabel{divisors-lemma-weakly-ass-zero@cref}{{[originallemma][0][]05AP}{[1][0][]source}{}{}{}}
\makeatother
\crefname{originallemma}{original lemma}{original lemmas}
\Crefname{originallemma}{Original lemma}{Original lemmas}
\makeatletter
\newlabel{lemma-bourbaki-finite-type-general-base}{{05IJ}{source}{}{sourcepointer-7}{}}
\newlabel{lemma-bourbaki-finite-type-general-base@cref}{{[originallemma][0][]05IJ}{[1][0][]source}{}{}{}}
\makeatother
\crefname{originallemma}{original lemma}{original lemmas}
\Crefname{originallemma}{Original lemma}{Original lemmas}
\makeatletter
\newlabel{lemma-finite-type-flat-at-point-local}{{05I7}{source}{}{sourcepointer-8}{}}
\newlabel{lemma-finite-type-flat-at-point-local@cref}{{[originallemma][0][]05I7}{[1][0][]source}{}{}{}}
\makeatother
\crefname{originallemma}{original lemma}{original lemmas}
\Crefname{originallemma}{Original lemma}{Original lemmas}
\makeatletter
\newlabel{modules-lemma-finite-type-surjective-on-stalk}{{01B8}{source}{}{sourcepointer-9}{}}
\newlabel{modules-lemma-finite-type-surjective-on-stalk@cref}{{[originallemma][0][]01B8}{[1][0][]source}{}{}{}}
\makeatother
\crefname{originallemma}{original lemma}{original lemmas}
\Crefname{originallemma}{Original lemma}{Original lemmas}
\makeatletter
\newlabel{more-morphisms-theorem-openness-flatness}{{0399}{source}{}{sourcepointer-10}{}}
\newlabel{more-morphisms-theorem-openness-flatness@cref}{{[originallemma][0][]0399}{[1][0][]source}{}{}{}}
\makeatother
\title{Openness of the flat locus for a finite type module over a locally finitely presented morphism without Noetherian hypotheses}
\author{The Stacks Project Authors}
\date{stacks-project:flat.tex:local-snapshot-20261011; GNU FDL 1.2 or later}
\begin{document}
\maketitle
\noindent\textit{Editorial note.} All delivered mathematical text is the author's own, copied from the pinned Stacks Project chapter file flat.tex; no statement or proof step was rewritten, reordered or removed.\par\medskip
\noindent\textit{Editorial note.} Preamble provenance: the chapter's own retained minimal preamble (sources/stacks-flat/retained-preamble.tex, lines 1-85) is a verbatim excerpt of the native Xy-pic setup and of the native theorem declarations and mathematical macros, followed by the project's own reference/slogan/history display-environment definitions copied mechanically from the shared preamble; the wrapper class and the amsmath/amssymb/amsthm/graphicx base packages are supplied by the assembler, and extra\_packages is left empty because everything needed is already in the retained preamble. No macro was invented or redefined.\par\medskip
\noindent\textit{Editorial note.} Same-chapter pointers are located at their \textbackslash{}label lines inside the pinned chapter file; cross-chapter pointers are located in the project's own tag list staged in the snapshot (sources/stacks-chow/tags.txt), at the verbatim lines carrying their canonical tags.\par\medskip
\noindent\textit{Allowed standard prerequisites (editorial).}\begin{itemize}
\item A finite type module flat at a point is finitely presented there, and a finitely presented module can be manufactured from the stalk. Same chapter flat.tex, label lemma-finite-type-flat-at-point-local (proof line 4121); Algebra, label algebra-lemma-construct-fp-module-from-stalk (proof line 4123). Produces a finitely presented module F' with a map to F that is an isomorphism at the chosen point, which is the first construction of the proof.
\item Openness of flatness and surjectivity of a map of finite type modules detected on a stalk. More on Morphisms, label more-morphisms-theorem-openness-flatness (proof line 4129); Modules, label modules-lemma-finite-type-surjective-on-stalk (proof line 4133); Algebra, label algebra-lemma-NAK (proof line 4135). Makes F' flat over the base after shrinking the source and the map from F' to F surjective after shrinking again, so that the proof is reduced to controlling the kernel.
\item Weakly associated points of a tensor product with a pulled back sheaf, finiteness of the associated points of a fibre, and weakly associated primes of a submodule. Same chapter flat.tex, label lemma-bourbaki-finite-type-general-base (proof line 4137); Divisors, labels divisors-lemma-finite-ass (proof line 4145) and divisors-lemma-ses-weakly-ass (proof line 4153); Algebra, labels algebra-lemma-silly (proof line 4147) and Divisors, label divisors-lemma-weakly-ass-zero (proof line 4158). The first bounds the weakly associated points of F' by the finitely many weakly associated points of the base times the associated points of the fibres, the second makes those finite, the third lets the weakly associated points be assumed to specialize into the chosen point, and the last two are used to conclude that a module whose weakly associated points are all generalizations of a point at which it vanishes is zero.
\end{itemize}\medskip
\ifdefined\hypertarget\hypertarget{sourcepointer-1}{}\fi
\noindent\textit{Source reference.} Original lemma 00DV: sources/stacks-chow/tags.txt, source lines 512--512. Navigation pointer for \textbackslash{}ref\{algebra-lemma-NAK\} used inside the counted span (line 4135); the argument that consumes it is retained in the span.\par\medskip
\ifdefined\hypertarget\hypertarget{sourcepointer-2}{}\fi
\noindent\textit{Source reference.} Original lemma 05GJ: sources/stacks-chow/tags.txt, source lines 7063--7063. Navigation pointer for \textbackslash{}ref\{algebra-lemma-construct-fp-module-from-stalk\} used inside the counted span (line 4123); the argument that consumes it is retained in the span.\par\medskip
\ifdefined\hypertarget\hypertarget{sourcepointer-3}{}\fi
\noindent\textit{Source reference.} Original lemma 00DS: sources/stacks-chow/tags.txt, source lines 508--508. Navigation pointer for \textbackslash{}ref\{algebra-lemma-silly\} used inside the counted span (line 4147); the argument that consumes it is retained in the span.\par\medskip
\ifdefined\hypertarget\hypertarget{sourcepointer-4}{}\fi
\noindent\textit{Source reference.} Original lemma 05AF: sources/stacks-chow/tags.txt, source lines 6848--6848. Navigation pointer for \textbackslash{}ref\{divisors-lemma-finite-ass\} used inside the counted span (line 4145); the argument that consumes it is retained in the span.\par\medskip
\ifdefined\hypertarget\hypertarget{sourcepointer-5}{}\fi
\noindent\textit{Source reference.} Original lemma 05AN: sources/stacks-chow/tags.txt, source lines 6856--6856. Navigation pointer for \textbackslash{}ref\{divisors-lemma-ses-weakly-ass\} used inside the counted span (line 4153); the argument that consumes it is retained in the span.\par\medskip
\ifdefined\hypertarget\hypertarget{sourcepointer-6}{}\fi
\noindent\textit{Source reference.} Original lemma 05AP: sources/stacks-chow/tags.txt, source lines 6857--6857. Navigation pointer for \textbackslash{}ref\{divisors-lemma-weakly-ass-zero\} used inside the counted span (line 4158); the argument that consumes it is retained in the span.\par\medskip
\ifdefined\hypertarget\hypertarget{sourcepointer-7}{}\fi
\noindent\textit{Source reference.} Original lemma 05IJ: sources/stacks-flat/pinned-original-source.tex, source lines 4079--4079. Navigation pointer for \textbackslash{}ref\{lemma-bourbaki-finite-type-general-base\} used inside the counted span (line 4137); the argument that consumes it is retained in the span.\par\medskip
\ifdefined\hypertarget\hypertarget{sourcepointer-8}{}\fi
\noindent\textit{Source reference.} Original lemma 05I7: sources/stacks-flat/pinned-original-source.tex, source lines 2879--2879. Navigation pointer for \textbackslash{}ref\{lemma-finite-type-flat-at-point-local\} used inside the counted span (line 4121); the argument that consumes it is retained in the span.\par\medskip
\ifdefined\hypertarget\hypertarget{sourcepointer-9}{}\fi
\noindent\textit{Source reference.} Original lemma 01B8: sources/stacks-chow/tags.txt, source lines 1717--1717. Navigation pointer for \textbackslash{}ref\{modules-lemma-finite-type-surjective-on-stalk\} used inside the counted span (line 4133); the argument that consumes it is retained in the span.\par\medskip
\ifdefined\hypertarget\hypertarget{sourcepointer-10}{}\fi
\noindent\textit{Source reference.} Original lemma 0399: sources/stacks-chow/tags.txt, source lines 4264--4264. Navigation pointer for \textbackslash{}ref\{more-morphisms-theorem-openness-flatness\} used inside the counted span (line 4129); the argument that consumes it is retained in the span.\par\medskip
\begin{theorem}
\label{theorem-finite-type-flat}
\begin{slogan}
The flat locus is open (non-Noetherian version).
\end{slogan}
Let $f : X \to S$ be a morphism of schemes.
Let $\mathcal{F}$ be a quasi-coherent $\mathcal{O}_X$-module.
Assume
\begin{enumerate}
\item $X \to S$ is locally of finite presentation,
\item $\mathcal{F}$ is an $\mathcal{O}_X$-module of finite type, and
\item the set of weakly associated points of $S$ is locally finite in $S$.
\end{enumerate}
Then $U = \{x \in X \mid \mathcal{F}\text{ flat at }x\text{ over }S\}$
is open in $X$ and $\mathcal{F}|_U$ is an $\mathcal{O}_U$-module
of finite presentation and flat over $S$.
\end{theorem}

\begin{proof}
Let $x \in X$ be such that $\mathcal{F}$ is flat at $x$ over $S$.
We have to find an open neighbourhood of $x$ such that $\mathcal{F}$ restricts
to a $S$-flat finitely presented module on this neighbourhood.
The problem is local on $X$ and $S$, hence we may assume that $X$ and $S$
are affine. As $\mathcal{F}_x$ is a finitely presented
$\mathcal{O}_{X, x}$-module by
Lemma \ref{lemma-finite-type-flat-at-point-local}
we conclude from
Algebra, Lemma \ref{algebra-lemma-construct-fp-module-from-stalk}
there exists a finitely presented $\mathcal{O}_X$-module $\mathcal{F}'$
and a map $\varphi : \mathcal{F}' \to \mathcal{F}$ which induces
an isomorphism $\varphi_x : \mathcal{F}'_x \to \mathcal{F}_x$. In particular
we see that $\mathcal{F}'$ is flat over $S$ at $x$, hence by openness
of flatness
More on Morphisms, Theorem \ref{more-morphisms-theorem-openness-flatness}
we see that after shrinking $X$ we may assume that
$\mathcal{F}'$ is flat over $S$. As $\mathcal{F}$ is of finite type
after shrinking $X$ we may assume that $\varphi$ is surjective, see
Modules, Lemma \ref{modules-lemma-finite-type-surjective-on-stalk}
or alternatively use Nakayama's lemma
(Algebra, Lemma \ref{algebra-lemma-NAK}).
By
Lemma \ref{lemma-bourbaki-finite-type-general-base}
we have
$$
\text{WeakAss}_X(\mathcal{F}') \subset
\bigcup\nolimits_{s \in \text{WeakAss}(S)} \text{Ass}_{X_s}(\mathcal{F}'_s)
$$
As $\text{WeakAss}(S)$ is finite by assumption and since
$\text{Ass}_{X_s}(\mathcal{F}'_s)$ is finite by
Divisors, Lemma \ref{divisors-lemma-finite-ass}
we conclude that $\text{WeakAss}_X(\mathcal{F}')$ is finite. Using
Algebra, Lemma \ref{algebra-lemma-silly}
we may, after shrinking $X$ once more, assume that
$\text{WeakAss}_X(\mathcal{F}')$ is contained in the generalization
of $x$. Now consider $\mathcal{K} = \Ker(\varphi)$. We have
$\text{WeakAss}_X(\mathcal{K}) \subset \text{WeakAss}_X(\mathcal{F}')$
(by
Divisors, Lemma \ref{divisors-lemma-ses-weakly-ass})
but on the other hand, $\varphi_x$ is an isomorphism, also $\varphi_{x'}$
is an isomorphism for all $x' \leadsto x$. We conclude that
$\text{WeakAss}_X(\mathcal{K}) = \emptyset$ whence
$\mathcal{K} = 0$ by
Divisors, Lemma \ref{divisors-lemma-weakly-ass-zero}.
\end{proof}

\end{document}
